% !TeX program = tectonic
\documentclass[11pt,a4paper]{article}
\usepackage{fontspec}
\usepackage[english]{babel}
\babelfont{rm}[Language=Default]{Liberation Serif}
\babelfont[Language=Default]{sf}{Carlito}
\babelfont[Language=Default]{tt}{DejaVu Sans Mono}
\usepackage{amsmath,amssymb,amsthm}
\usepackage{geometry}
\geometry{a4paper, top=2.4cm, bottom=2.4cm, left=2.4cm, right=2.4cm}
\usepackage{booktabs,tabularx,enumitem,xcolor,tcolorbox}
\tcbuselibrary{breakable,skins}
\definecolor{linkblue}{HTML}{1F4E79}
\definecolor{statgreen}{HTML}{2F6B3A}
\definecolor{goldacc}{HTML}{8B7E5A}
\usepackage[colorlinks=true,linkcolor=linkblue,urlcolor=linkblue,unicode=true]{hyperref}
\theoremstyle{plain}
\newtheorem{theorem}{Theorem}
\newtheorem{lemma}[theorem]{Lemma}
\newtheorem{proposition}[theorem]{Proposition}
\newtheorem{corollary}[theorem]{Corollary}
\theoremstyle{definition}
\newtheorem{definition}[theorem]{Definition}
\theoremstyle{remark}
\newtheorem{remark}[theorem]{Remark}
\newtcolorbox{statusbox}[1][]{colback=statgreen!5,colframe=statgreen!75,
  title={\textbf{Summary of results:} #1},fonttitle=\small,boxrule=0.7pt,arc=2pt,breakable}
\newtcolorbox{databox}[1][]{colback=goldacc!6,colframe=goldacc!80,
  title={\textbf{Protocol data:} #1},fonttitle=\small,boxrule=0.7pt,arc=2pt,breakable}
\newtcolorbox{verifybox}[1][]{colback=linkblue!5,colframe=linkblue!80,
  title={\textbf{Verification:} #1},fonttitle=\small,boxrule=0.7pt,arc=2pt,breakable}
\widowpenalty=10000
\clubpenalty=10000
\setlength{\parskip}{0.35em}
\newcommand{\CC}{\mathbb{C}}
\newcommand{\RR}{\mathbb{R}}
\newcommand{\ZZ}{\mathbb{Z}}
\newcommand{\QQ}{\mathbb{Q}}
\newcommand{\Ga}{\Gamma}
\newcommand{\im}{\operatorname{Im}}
\newcommand{\re}{\operatorname{Re}}
\newcommand{\disc}{\operatorname{disc}}
\begin{document}
\begin{center}
{\LARGE\bfseries Certificate I: The Hurwitz Rung $N=7$}\\[0.4em]
{\large The cyclic cubic of $2\cos(2\pi/7)$, the Cardano pairing, and the $\pi/7$ phase}\\[0.6em]
{\normalsize Isaev Iskhak Khamzatovich}\\[0.2em]
{\small The ``Dynamic Principle'' program --- the \texttt{hodge-laboratory} repository}\\[0.15em]
{\small Standalone theorem monograph · Monograph 17/20}
\end{center}
\vspace{0.4em}
{\color{linkblue}\hrule height 0.8pt}
\vspace{0.8em}

\section{Setting}

With the program extended to the four rungs $N=7,9,15,30$
(roadmap v1.3), the ladder acquired two cubic foundations: the
Hurwitz rung $N=7$ and the Macbeath rung $N=9$. Certificate~I is
the first of the two new anchors: it secures, at level $N=7$, the
entire pipeline of the program --- the conductor census of
eigencharacters, the exact cubic layer of $b_{\mathrm{Ch}}(7)$ in
$\ZZ[\zeta]/(\Phi_7)$, and the Cardano closed form in the
$\mathrm{casus\ irreducibilis}$ regime.

The rung is named after Adolf Hurwitz: the polynomial
$x^3+x^2-2x-1$ is the minimal polynomial of $2\cos(2\pi/7)$, a
generator of the maximal totally real subfield of the cyclotomic
field $\QQ(\zeta_7)$ --- the very field that lives in the
discriminant $\Delta=-343=-7^3$ of the Klein quartic (monograph
16/20). Note, however, that the $N=7$ stand certifies not the
Klein quartic itself (which has its own ``Klein'' stand and
monograph 16/20) but the Fermat-level pipeline at $N=7$: the
Fermat curve $X_7$ has genus $15$, and every mechanism of the
program --- census, periods, $\Ga$-monomials --- operates at this
level unchanged.

\begin{theorem}[I1 --- the census of the prime level]\label{th:I1}
For the Fermat curve $X_7$ every eigencharacter $(a,b)$ with
$1\le a,b$, $a+b\le 6$ has conductor $d=7$; the conductor census
consists of a single class, $\{7\colon 15\}$, and
$\sum_d h_d = 15 = (7-1)(7-2)/2 = g$. Both census schemes (the
direct one and the M\"obius one) produce the same histogram.
\end{theorem}

\begin{proof}
The level $N=7$ is prime, so for any eigencharacter $(a,b)$ with
$1\le a,b$, $a+b\le 6$ we have $a,b\not\equiv 0 \pmod 7$, hence
$\gcd(7,a,b)=1$ and the conductor is $d = 7/\gcd(7,a,b) = 7$ ---
the only one. The number of such pairs $(a,b)$ is the count of
lattice points in the triangle $1\le a,b$, $a+b\le 6$: for each
$a$ from $1$ to $5$ there are exactly $6-a$ values of $b$, giving
$\sum_{a=1}^{5}(6-a)=5+4+3+2+1=15$. The M\"obius scheme:
$h_7 = c(7)-c(1) = 15-0 = 15$, where
$c(k)=(k-1)(k-2)/2$ is the number of eigencharacters of level
$k$. Both schemes are integer identities; their agreement is
deterministic.
\end{proof}

\begin{theorem}[I2 --- the cubic layer of $b_{\mathrm{Ch}}(7)$]\label{th:I2}
The numbers $x_1=\zeta+\zeta^{-1}$, $x_2=\zeta^2+\zeta^{-2}$,
$x_3=\zeta^3+\zeta^{-3}$ with $\zeta=e^{2\pi i/7}$ have exact
elementary symmetric polynomials
\[
  s_1=x_1+x_2+x_3=-1,\qquad
  s_2=x_1x_2+x_1x_3+x_2x_3=-2,\qquad
  s_3=x_1x_2x_3=1,
\]
computed by pure integer arithmetic in $\ZZ[\zeta]/(\Phi_7)$;
therefore $x_1$ is a root of the cubic $x^3+x^2-2x-1$, and this
is its exact minimal polynomial over $\QQ$.
\end{theorem}

\begin{proof}
\emph{Sum.} $s_1=\zeta+\zeta^6+\zeta^2+\zeta^5+\zeta^3+\zeta^4$
is the sum of all nontrivial powers of $\zeta$, i.e.
$s_1=-1$ (the $\Phi_7$-sum, since $1+\zeta+\dots+\zeta^6=0$).

\emph{Pairwise products.} Expand in $\ZZ[\zeta]/(\Phi_7)$ using
$\zeta^{-k}=\zeta^{7-k}$:
$x_1x_2=(\zeta+\zeta^6)(\zeta^2+\zeta^5)=\zeta^3+\zeta^6+\zeta+\zeta^4$;
$x_1x_3=(\zeta+\zeta^6)(\zeta^3+\zeta^4)=\zeta^4+\zeta^5+\zeta^2+\zeta^3$;
$x_2x_3=(\zeta^2+\zeta^5)(\zeta^3+\zeta^4)=\zeta^5+\zeta^6+\zeta+\zeta^2$.
Adding, each of the six powers $\zeta^1,\dots,\zeta^6$ occurs
exactly twice:
$s_2=2(\zeta+\zeta^2+\dots+\zeta^6)=2\cdot(-1)=-2$.

\emph{Product.}
$x_1x_2x_3=(\zeta^3+\zeta^6+\zeta+\zeta^4)(\zeta^3+\zeta^4)
=\zeta^6+\zeta^7+\zeta^9+\zeta^{10}+\zeta^4+\zeta^5+\zeta^7+\zeta^8
=\zeta^6+1+\zeta^2+\zeta^3+\zeta^4+\zeta^5+1+\zeta
=2+(\zeta+\dots+\zeta^6)=2-1=1$.

By Vieta's formulas the cubic with roots $x_1,x_2,x_3$ is
$x^3-s_1x^2+s_2x-s_3=x^3+x^2-2x-1$. Minimality: the Galois group
$\mathrm{Gal}(\QQ(\zeta_7)^+/\QQ)\cong(\ZZ/7)^\times/\{\pm1\}\cong C_3$
acts by multiplication of exponents by $2$:
$1\mapsto2\mapsto4\equiv-3\mapsto-1\equiv1$ --- a cycle on the
three classes $\{1,6\},\{2,5\},\{3,4\}$; the action is
transitive, so $[\QQ(x_1):\QQ]=3$. Fingerprint: modulo $2$ the
cubic $x^3+x^2-2x-1\equiv x^3+x^2+1$ has no root in $\mathbf{F}_2$
($0\mapsto1$, $1\mapsto1$) --- it is irreducible, independently
confirming the degree $3$.
\end{proof}

\begin{theorem}[I3 --- the Cardano form and the pairing]\label{th:I3}
In terms of $t=\frac{7}{54}+\frac{7\sqrt{-3}}{18}$ and
$\bar t=\frac{7}{54}-\frac{7\sqrt{-3}}{18}$ the largest root of
the cubic is
\[
  2\cos\frac{2\pi}{7}=\sqrt[3]{t}+\sqrt[3]{\bar t}-\frac13,
  \qquad
  b_{\mathrm{Ch}}(7)=1-\cos\frac{2\pi}{7}
  =\frac76-\frac{\sqrt[3]{t}+\sqrt[3]{\bar t}}{2},
\]
and the moduli of the Cardano arguments are paired by the exact
identity $|t|^2=(7/9)^3$, so that the product of the cube roots
is $\sqrt[3]{t}\cdot\sqrt[3]{\bar t}=7/9$ exactly.
\end{theorem}

\begin{proof}
The substitution $x=y-\frac13$ removes the quadratic term:
$p=-\frac73$, $q=-\frac{7}{27}$, and the cubic becomes
$y^3-\frac73y-\frac7{27}=0$. The Cardano discriminant is
\[
  \Delta=\frac{q^2}{4}+\frac{p^3}{27}
  =\frac{49}{2916}-\frac{343}{729}
  =\frac{49-1372}{2916}=-\frac{1323}{2916}=-\frac{49}{108}<0
\]
--- the classical \emph{casus irreducibilis}: all three roots are
real, yet expressible only through complex cube roots. Hence
$\sqrt{\Delta}=\frac{7i\sqrt3}{18}$ and
\[
  y=\sqrt[3]{-\frac q2+\sqrt\Delta}+\sqrt[3]{-\frac q2-\sqrt\Delta}
  =\sqrt[3]{\frac{7}{54}+\frac{7\sqrt{-3}}{18}}
  +\sqrt[3]{\frac{7}{54}-\frac{7\sqrt{-3}}{18}},
\]
so $x_1=y-\frac13$ and $b_{\mathrm{Ch}}(7)=1-\frac{x_1}{2}
=\frac76-\frac{\sqrt[3]{t}+\sqrt[3]{\bar t}}{2}$.

\emph{The pairing.} $|t|^2=\frac{49}{2916}+\frac{147}{324}
=\frac{49+1323}{2916}=\frac{1372}{2916}=\frac{343}{729}=\left(\frac79\right)^3$;
therefore $\sqrt[3]{t}\cdot\sqrt[3]{\bar t}=\sqrt[3]{|t|^2}
=\sqrt[3]{(7/9)^3}=7/9$ --- an exact integer identity (the
``the Cardano pairing'' comment in the laboratory code).

\emph{Vi\`ete form.} $|t|^{1/3}=\sqrt{7}/3$ and
$\arg t=\arctan(3\sqrt3)$, so the principal cube root is
$u=\sqrt[3]{t}=\frac{\sqrt7}{3}e^{i\theta}$ with
$\theta=\frac13\arctan(3\sqrt3)$, and
$2\cos\frac{2\pi}{7}=2\re u-\frac13
=\frac{2\sqrt7}{3}\cos\!\Big(\frac13\arctan(3\sqrt3)\Big)-\frac13$
--- the real trigonometric form (the trisection of the angle
$\arctan(3\sqrt3)\approx79.1^\circ$).

\emph{Isolation.} The roots $x_1\approx2.24698$,
$x_2\approx-0.55496$, $x_3\approx-1.69202$ are strictly ordered;
the numeric check at working precision (dps 60) gives a relative
deviation of the Cardano form from $1-\cos(2\pi/7)$ equal to
$2.1\cdot10^{-61}$ --- agreement within working precision,
pinning the largest root.
\end{proof}

\begin{remark}
The phase $\pi/7$ anchors the quantum ladder of the program: item
G6 of certificate G secures the phases $\pi/7,\pi/15,\pi/30$ as
levels of the single construction $\delta=\pi/N$, and the isotypy
$(1,1,0,1)$ of the Klein quartic ties $\zeta_7$ to its periods.
The rung $N=7$ is thus the arithmetic foundation of the whole
heptadic branch.
\end{remark}

\begin{databox}[Certificate I --- the stand numbers]
Genus $g=15$; census $\{7\colon 15\}$ (both schemes); minimal
cubic $x^3+x^2-2x-1$; Vieta $(s_1,s_2,s_3)=(-1,-2,1)$ exactly in
$\ZZ[\zeta]/(\Phi_7)$; the cubic is irreducible over
$\mathrm{GF}(2)$; Cardano form
$b_{\mathrm{Ch}}(7)=\frac76-\frac{\sqrt[3]{t}+\sqrt[3]{\bar t}}{2}$,
$t=\frac{7}{54}\pm\frac{7\sqrt{-3}}{18}$; the pairing
$\sqrt[3]{t}\sqrt[3]{\bar t}=7/9$;
$b_{\mathrm{Ch}}(7)=0.376510198141266469474995115995760\ldots$;
protocol relative errors: Cardano $2.0\cdot10^{-36}$, periods
(spot sample) $1.0\cdot10^{-36}$ (dps 35); phase $\pi/7$.
\end{databox}

\begin{statusbox}[summary]
Certificate I accepted ($5/5$): the census of the prime level by
two schemes, the exact integer Vieta data in
$\ZZ[\zeta]/(\Phi_7)$, irreducibility over $\mathrm{GF}(2)$, the
Cardano form in the $\mathrm{casus\ irreducibilis}$ with the exact
pairing $7/9$, and the phase $\pi/7$. The cubic foundation of the
ladder is established.
\end{statusbox}

\begin{verifybox}[how to reproduce]
\texttt{python3 laboratory.py} $\to$ ``Certificates $\to$ I'';
``Stands $\to$ N=7''; batch mode: a scenario with the
\texttt{bch} item ($N=7$); \texttt{--check-baseline} cross-checks
$g=15$, the symmetric data $(-1,-2,1)$ and the errors against the
reference.
\end{verifybox}

\vfill
{\color{linkblue}\hrule height 0.6pt}
\vspace{0.5em}
{\small\color{gray}
License: individual exclusive license of Isaev Iskhak Khamzatovich.
All rights to the computations, formulas and texts belong to the author.
Verification: \texttt{python3 laboratory.py} (``Stands''/``Certificates'').
}
\end{document}
